\documentclass[12pt,a4paper]{amsart}
\usepackage{amssymb,mathtools}
\usepackage[margin=1in]{geometry}
\usepackage[colorlinks=true,urlcolor=blue,linkcolor=blue,citecolor=blue]{hyperref}
\newtheorem*{theorem}{Theorem}
\theoremstyle{definition}
\newtheorem*{setup}{Setup}
\theoremstyle{plain}
\newtheorem{lemma}{Lemma}[section]
\newtheorem{proposition}[lemma]{Proposition}
\newtheorem{corollary}[lemma]{Corollary}
\newcommand{\Q}{\mathbb Q}
\newcommand{\C}{\mathbb C}
\newcommand{\OO}{\mathcal O}
\newcommand{\mult}{\operatorname{mult}}
\newcommand{\nef}{\operatorname{n}}
\title{Nef Abundance: Threefold Behavior}
\date{October 10, 2026}
\begin{document}
\maketitle


\section{Theorem and proof}\label{sec:hard-case}

\begin{theorem}\label{thm:full-nef-nonvanishing}
Let $X$ be a smooth projective threefold over $\C$ such that
$K_X\sim0$ and $q(X)=0$. Let $D$ be a nef Cartier divisor on $X$
such that $\nef(X,D)=3$ and $D^3=0$.
Then
\[
 H^0(X,\OO_X(mD))\ne0
 \qquad\text{for some integer }m>0.
\]
\end{theorem}

\begin{proof}
Suppose for contradiction that
\begin{equation}\label{eq:counterexample-assumption}
 \kappa(D)=-\infty,\qquad\text{equivalently }
 H^0(X,\OO_X(mD))=0\text{ for every integer }m>0.
\end{equation}

\begin{setup}
\leavevmode
\begin{itemize}
\setlength{\itemsep}{2pt}
\setlength{\parsep}{0pt}
\item Fix an ample Cartier divisor $A$.
\item For integers $n\ge1$, put
\begin{equation}\label{eq:Pt-setting}
 P_n=r_n(D+A/n),
\end{equation}
The positive integers $r_n$ are chosen in the next item. Each $P_n$
is ample and rational Cartier; whenever $|kP_n|$ is used, $k$ is a
positive integer such that $kP_n$ is Cartier.
\item Fix $\epsilon\in\Q_{>0}$ and choose $r_n$ to be the nearest
positive integer to $\epsilon/((D+A/n)^3)^{1/3}$. Since $D^3=0$,
Lemma~\ref{lem:hd-normalization} gives
$r_n\to\infty$ and $P_n^3\to\epsilon^3$.
\item Let $D_i$ and $P_{n,i}$ denote pullbacks from the two factors
of $X\times X$, and then to
\[
 Z=\mathbb P_{X\times X}
       \bigl(\OO_{X\times X}(D_1)\oplus\OO_{X\times X}(D_2)\bigr),
 \qquad \xi=c_1(\OO_Z(1)).
\]
For a positive integer $Q$, put $R_n=P_{n,1}+P_{n,2}+Q\xi$.
The torus open is the complement of the two summand sections.
\end{itemize}
We take $\epsilon$ sufficiently small that
\[
 (14\epsilon)^3<\tfrac14,\qquad80\epsilon<\tfrac34.
\]

\end{setup}

Under the preceding assumptions, we will prove the following three
lemmas in Sections~\ref{sec:geometric-inputs}, \ref{sec:large-jets},
and \ref{sec:frobenius}, respectively.

\begin{lemma}\label{lem:key-scalar}
For every fixed integer $Q>0$ and all sufficiently large $n$,
at a very general point $x\in X$, the blowup
$\pi_x:\widehat X\to X$, with exceptional divisor $J$, admits a
class $\gamma=f_*(A_1A_2)$ such that
\[
 J\gamma>0,\qquad
 \pi_x^*(P_n+\lambda D)\gamma\le40\epsilon J\gamma
 \quad\text{for all }\lambda\in[0,Q].
\]
Here $f:Y\to\widehat X$ is a birational morphism from a smooth
projective variety, and $A_1,A_2$ are ample Cartier divisors on $Y$.
\end{lemma}

\begin{lemma}\label{lem:key-jets}
There exists an integer $Q>0$ such that, for all sufficiently large
$n$, there exist a point $z$ in the torus open of $Z$ and an integer
$k_0>0$ with $k_0P_n$ Cartier for which the evaluation map
\[
 H^0(Z,\OO_Z(k_0R_n))\longrightarrow
 \frac{\OO_{Z,z}(k_0R_n)}
 {\mathfrak m_z^{8k_0+1}\OO_{Z,z}(k_0R_n)}
\]
is surjective.
\end{lemma}

\begin{lemma}\label{lem:key-contradiction}
Fix $Q$ as in Lemma~\ref{lem:key-jets} and $n$ with
$r_n>Q+1$ and $P_n^3<(2\epsilon)^3$. Put $P=P_n$.
Fix the data in Lemmas~\ref{lem:key-scalar} and \ref{lem:key-jets},
and a smooth complete intersection flag of successive classes
$l_1P,l_2P$ whose final curve $C$ satisfies
$\mu_{\min}(\Omega_X^1|_C)\ge0$.
On smooth reductions $Y$ in arbitrarily large characteristics $p$,
there exists a nonzero effective Cartier divisor $E_p$ on
$Y'\times Y'$ such that
\[
 \tfrac34p^6<\mult_{(x',x')}(E_p)
 \le\bigl(80\epsilon+O(p^{-1})\bigr)p^6.
\]
Here $Y'$ is the relative Frobenius twist of $Y$, and $x'$ is
the point induced by $x$; the implied constant is independent of $p$.
\end{lemma}

We now complete the proof.

By Lemma~\ref{lem:key-jets}, choose $Q$ independently of $n$.
Then choose $n$ sufficiently large that Lemmas~\ref{lem:key-scalar}
and \ref{lem:key-jets} apply, $r_n>Q+1$, and
$P_n^3<(2\epsilon)^3$. Set $P=P_n$ and fix a jet system supplied
by Lemma~\ref{lem:key-jets}.

Choose a very general point $x\in X$. Lemma~\ref{lem:key-scalar}
supplies a strongly movable witness on the blowup at $x$.
By Lemma~\ref{lem:cotangent-flag}, choose a smooth complete
intersection flag of successive classes $l_1P,l_2P$, whose final
curve $C$ satisfies $\mu_{\min}(\Omega_X^1|_C)\ge0$.
Fix the jet system, witness, and flag over $\C$ before reduction.

Dividing the inequalities in Lemma~\ref{lem:key-contradiction}
by $p^6$ gives
\[
 \tfrac34<80\epsilon+O(p^{-1})
\]
for arbitrarily large $p$, contradicting $80\epsilon<3/4$.
Thus \eqref{eq:counterexample-assumption} is false, and
$H^0(X,\OO_X(mD))\ne0$ for some integer $m>0$.\qedhere
\end{proof}

\section{Proof of Lemma~\ref{lem:key-scalar}}\label{sec:geometric-inputs}

Write $\nu(D)=\max\{j:D^jA^{3-j}>0\}\in\{1,2\}$.
We first bound multiplicities and then obtain the required movable
class by duality. The same base-component estimates are used in
Section~\ref{sec:large-jets}.
All auxiliary varieties in these two sections are over $\C$.

\subsection{Normalization and numerical inputs}\label{sec:lower-dimensional-inputs}

\begin{lemma}\label{lem:hd-normalization}
For every $\epsilon>0$ and
as $n\to\infty$ through the positive integers, there are
positive integers $r_n\to\infty$ such that
\[
 P_n=r_n(D+A/n),\qquad P_n^3\longrightarrow\epsilon^3.
\]
The divisor $P_n$ is ample and rational Cartier; its Cartier multiples
are chosen after this normalization.
\end{lemma}
\begin{proof}
Put $a_n=((D+A/n)^3)^{1/3}$. Then $a_n>0$ and $a_n\to0$.
Choose $r_n$ to be the nearest positive integer to $\epsilon/a_n$.
For sufficiently large $n$,
$|r_na_n-\epsilon|\le a_n/2$, which proves the assertions.
\end{proof}

\begin{lemma}\label{lem:lower-dimensional}
Let $(Y,B)$ be an auxiliary projective klt pair of dimension
$f\in\{1,2\}$, with $T=K_Y+B$ pseudo-effective, and let $M$ be
nef and rational Cartier. Then $T+M$ has an effective rational
numerical representative. Moreover, if for some $a>0$
\[
 (T+aM)\cdot C>0
\]
for every integral curve through a sufficiently very general point,
then $T+cM$ is big for every rational $c>0$.
\end{lemma}
\begin{proof}
Numerical nonvanishing on surfaces is
\cite[Corollary~C]{LP}; on curves it follows from the degree criterion. For the second assertion, the curve case is immediate.
In dimension two run the ordinary $(K_Y+B)$-minimal model program,
after a small $\Q$-factorialization. Its output is a morphism
$g:Y\to Y_0$ with $U=K_{Y_0}+B_0$ nef. The pushforward
$N=g_*M$ is nef. Negativity gives
\[
 T=g^*U+E_T,\qquad M=g^*N-E_M,
\]
where $E_T,E_M\ge0$ are exceptional; $E_T$ has positive
coefficient on every contracted divisor. Thus, for small rational
$s>0$, $E_T-sE_M\ge0$.

By \cite[Corollary~C]{LP}, $U+sN$ is numerically equivalent
to a semiample rational divisor. If it were not big, the resulting
morphism would have positive-dimensional fibres. A general fibre,
or a general curve in such a fibre when the image is a point,
avoids the finitely many points $g(\operatorname{Exc}g)$.
Since $U$ and $N$ are nef, this curve has both $U$- and $N$-degree
zero. Its inverse image contradicts the stated full positivity of
$T+aM$. Consequently $U+sN$, and hence
$T+sM=g^*(U+sN)+E_T-sE_M$, is big.
For $c>0$, choose $0<\lambda<\min\{1,c/s\}$ and write
\[
 T+cM=\lambda(T+sM)+(1-\lambda)T+(c-\lambda s)M.
\]
The last two summands are pseudo-effective, so the sum is big.
This proves exactly the dimension-one and dimension-two inputs
used in \cite[Sections~3--5]{NS}, without assuming threefold
nef nonvanishing.
\end{proof}

\subsection{Moving base components}

\begin{lemma}\label{lem:hd-base-degree}
Let $W$ be smooth projective of dimension $d$, let $H$ be a nef
Cartier divisor on $W$, and let
$\Lambda\subseteq|H|$ be a nonempty linear system
with base ideal $\mathfrak b$. Let $V$ be an integral component of
its base locus, of codimension $c\ge1$. If
\[
 \mathfrak b\OO_{W,\eta_V}\subseteq
 \mathfrak m_{\eta_V}^{h}\qquad(h\in\mathbb Z_{>0}),
\]
then
\[
 h^c\,H^{d-c}\cdot V\le H^d.
\]
In particular, if $x$ is an isolated base point and every member
has multiplicity at least a real number $r\ge0$ at $x$, then
$r^d\le H^d$.
\end{lemma}
\begin{proof}
Intersect successively with general members of $\Lambda$.
Before each next cut, remove the components already contained in
the base locus. Each remaining component is cut properly, and the
intersection cycle is effective. Its $H$-degree equals the $H$-degree
of the preceding cycle; discarding effective components can only
decrease this degree, since $H$ is nef. Starting with $[W]$, the
$H$-degree of the codimension-$c$ cycle is therefore at most $H^d$.

At $\eta_V$, the localized base ideal is primary for the maximal
ideal. No discarded component of smaller codimension contains
$\eta_V$, because $V$ is an irreducible component of the base locus.
The first $c$ general equations $s_1,\ldots,s_c$ form a parameter
ideal $I$ in the regular local ring $R=\OO_{W,\eta_V}$.
The coefficient of $V$ in the intersection cycle is
$\operatorname{length}(R/I)$. Since $I\subseteq\mathfrak m^h$,
monotonicity of Hilbert--Samuel multiplicity gives
\[
 \operatorname{length}(R/I)=e(I,R)
 \ge e(\mathfrak m^h,R)=h^c.
\]
Here the first equality uses that a parameter sequence in $R$ is
regular, and the last equality uses regularity of $R$.
This proves the inequality. For an isolated point take $c=d$ and
$h=\max\{1,\lceil r\rceil\}$.
\end{proof}

\begin{lemma}\label{lem:finite-subadjunction-model}
Let $W$ be smooth projective, $\Theta\ge0$ rational, and $V\subset W$
integral of positive dimension. Suppose $(W,\Theta)$ is lc at
$\eta_V$ and has an lc place centred on $V$.
For $L=(K_W+\Theta)|_{V^\nu}$ there exist a finite surjective
morphism $\eta:V'\to V^\nu$, with $V'$ normal, and rational
divisors $B_{V'}\ge0$ and $M_{V'}$, the trace of a b-nef divisor,
such that
\[
 \eta^*L\sim_\Q K_{V'}+B_{V'}+M_{V'}.
\]
\end{lemma}
\begin{proof}
We use the construction in \cite[Lemma~7.1]{LA}.
The dlt blowup \cite[Theorem~2.10]{FH} gives
\[
 p^*(K_W+\Theta)=K_Q+B_Q+E_Q,
 \qquad (Q,B_Q)\text{ dlt},\quad E_Q\ge0,
\]
where $B_Q$ is the truncated boundary and the effective surplus
$E_Q$ is absent near $\eta_V$. Choose a dlt stratum $S$ minimal
among those dominating $V$. Adjunction, followed by adding
$E_Q|_S$, gives an effective rational divisor $B_S$ and
\[
 K_S+B_S\sim_\Q(p|_S)^*((K_W+\Theta)|_V).
\]
The stratum is not contained in $E_Q$. Minimality makes
$(S,B_S)$ klt over $\eta_V$, although it need not be globally lc.

Factor through the normalization and the connected-fibre base:
$S\xrightarrow{q}V'\xrightarrow{\eta}V^\nu\to V$, with $\eta$
finite. The discrepancy-round-up rank condition for $q$ is one:
over the generic point the effective klt boundary contributes
zero round-up on strict divisors and allows poles only along
exceptional divisors. Its pushforward to normal $S$ is therefore
$\OO_S$, and connectedness gives rank one. Thus this is a
klt-trivial fibration in \cite[Definition~3.1]{FG}. Its formula is
\[
 \eta^*((K_W+\Theta)|_{V^\nu})
 \sim_\Q K_{V'}+B_{V'}+M_{V'},
\]
with b-nef moduli divisor $M$ by \cite[Theorem~3.5]{FG}.
The discriminant on this initial base is effective: over a prime
divisor, an inverse-image component of multiplicity $m\ge1$ and
boundary coefficient $b\ge0$ bounds its lc threshold by
$(1-b)/m\le1$. A negative threshold still gives a positive
discriminant coefficient. Only generic sub-klt singularities are
required for this use of the formula.
\end{proof}

\begin{lemma}\label{lem:generic-subadjunction}
Under the hypotheses of Lemma~\ref{lem:finite-subadjunction-model},
let $f=\dim V$, let $P_V$ be nef and rational Cartier, and let
$\rho:Y\to V$ be a resolution. Then
\[
 K_Y(\rho^*P_V)^{f-1}\le(K_W+\Theta)|_V P_V^{f-1}.
\]
\end{lemma}
\begin{proof}
Apply Lemma~\ref{lem:finite-subadjunction-model} and intersect with
$f-1$ copies of the pulled-back $P_V$. The effective discriminant
and the moduli trace have nonnegative intersection; for the latter,
push down from a higher base where the moduli divisor is nef.
For compatible canonical Weil cycles,
\[
 \eta_*K_{V'}=(\deg\eta)K_{V^\nu}+\eta_*R_\eta,
 \qquad R_\eta\ge0.
\]
Discard ramification and divide by $\deg\eta$. Pushing down $K_Y$
to $V^\nu$ changes it only by exceptional cycles, whose testing
intersection is zero. No rational Cartier assumption on
$K_{V^\nu}$ is needed.
\end{proof}

\begin{lemma}\label{lem:canonical-centre}
Let $W$ be smooth projective of any dimension, and let $P_W$
be nef and rational Cartier. Suppose a positive-dimensional integral base component $V$
of a subseries of $|kP_W|$ is isolated at its generic point,
of codimension $a$, and its base ideal there is contained in
$\mathfrak m^h$. Then on a resolution $\rho:Y\to V$,
\[
 K_Y(\rho^*P_W)^{\dim V-1}
 \le (K_W+cP_W)|_V P_W^{\dim V-1},
 \qquad 0<c\le ak/h.
\]
\end{lemma}
\begin{proof}
The generic log canonical threshold of the base ideal satisfies
$\lambda\le a/h$, as seen by blowing up the smooth generic
centre. On a log resolution, average sufficiently many general
members of the subseries with total coefficient $\lambda$.
This gives an actual effective boundary
$\Theta\sim_{\Q}k\lambda P_W$ which is lc at the generic
point of $V$ and has an lc place centred on $V$.
Lemma~\ref{lem:generic-subadjunction} gives the inequality with
$c=k\lambda$. The construction uses actual general members of
the given subseries, and requires lc only at $\eta_V$.
Both the threshold bound and subadjunction are dimension-independent.
\end{proof}

\begin{lemma}\label{lem:hd-dense-multiplicity}
Let $H$ be a Cartier divisor on $X$, let
$\Lambda\subseteq|H|$ be a fixed nonempty linear system, and let
$r\ge0$ be real. If a Zariski dense set of points $x$ admits
$E_x\in\Lambda$ with $\mult_x(E_x)>r$, then every point does.
\end{lemma}
\begin{proof}
Use local parameters $u,v,w$ and homogeneous parameters
$[t_0:\cdots:t_m]$ on the fixed linear system. Write the universal
equation as $f=\sum_i t_i f_i$, and denote its coefficient functions
by $c_{i,\alpha}(x)$, where $\alpha\in\mathbb Z_{\ge0}^3$.
The incidence $Z$ is defined locally by
\[
 \sum_i t_i c_{i,\alpha}(x)=0\qquad(|\alpha|\le\lfloor r\rfloor).
\]
These coefficients are the finite local jet coefficients. Thus
$Z$ is closed, and its proper projection has closed image.
An image containing a Zariski dense subset equals $X$.

\end{proof}

\begin{lemma}\label{lem:jet-normal-containment}
Let $W$ be smooth projective, $L$ a line bundle, and $z:T\to W$
a morphism from a smooth integral variety. Fix integers
$m_1>m_0\ge0$ and $1\le h\le m_1-m_0+1$.
Suppose the spaces
\[
 \mathcal H_i(t)=H^0(W,L\otimes\mathfrak m_{z(t)}^{m_i})
 \quad(i=0,1)
\]
form vector bundles, and a reduced family $\Gamma\subset T\times W$
dominates $W$, with each general fibre $\Gamma_t$ a component of
both base loci. Then the higher base ideal satisfies
$\mathfrak b_{1,t}\subset I_{\Gamma_t}^{h}$ at a general incidence point.
\end{lemma}
\begin{proof}
A frame of the higher kernel has the form
$s(t,w)=\sum_\alpha a_\alpha(t)s_\alpha(w)$ for fixed global sections.
A parameter derivative of order $j<h$ loses at most $j$ orders
at the moving mark, so belongs to the full lower kernel and
vanishes on $\Gamma$.
At a general smooth point, dominance of $\Gamma\to W$ makes
parameter directions span the normal space. Choose transverse
parameters $t'$ and write $\Gamma$ locally as
$u=t'-F(t'',w)=0$. Differentiation in $t'$ with $t'',w$ fixed
is differentiation in $u$. Thus all terms of degree less than $h$
in the normal parameters vanish, giving $s\in I_\Gamma^h$.
Characteristic zero and faithful flatness of the analytic local
ring give the algebraic containment. Restricting the finite frame
to a general parameter fibre proves the assertion, as in
\cite[Lemma~7.3]{LA}.
\end{proof}

\begin{lemma}\label{lem:restricted-components}
Let $g:W\to B$ be a morphism from a smooth projective variety,
and let an integral component $V$ of the base locus of a linear
system have base ideal $\mathfrak b\subset I_V^h$.
Let $F$ be the reduced component through a general point of a
fibre of $V\to g(V)$. Assume the ambient fibre $W_b$ is smooth
and integral and $F\subsetneq W_b$. Then the restricted system
is nonzero, $F$ is a base component, and
\[
 \mathfrak b|_{W_b}\subset I_F^h
\]
near $\eta_F$.
\end{lemma}
\begin{proof}
Remove the other base components and work where $V\to g(V)$ is
smooth over its actual image. Its scheme fibre is reduced there,
hence $I_V\OO_{W_b}=I_F$ and $I_V^h\OO_{W_b}=I_F^h$.
The local base support is exactly $F$.
Since $F$ is proper, not all restricted sections vanish identically
on $W_b$. The resulting base ideal is primary at $\eta_F$.
This is the actual-image restriction in \cite[Lemma~7.4]{LA}.
\end{proof}

\begin{lemma}\label{lem:full-kernels}
Let $W$ be a smooth projective variety of dimension $d$, and let
$P_W$ be ample and rational Cartier. Fix $kP_W$ Cartier and levels
$\tau_0<\cdots<\tau_d$ with common gap $\delta>0$ and
$k\delta\ge4$. Suppose that, at a very general marked point $z$,
all full spaces
\[
 \mathcal H_i(z)=H^0(W,\OO_W(kP_W)\otimes\mathfrak m_z^{\lceil k\tau_i\rceil})
\]
are nonzero, and the lowest base locus has a positive-dimensional
component through $z$. Then there is an algebraic covering family
of proper integral components $V_t$, common to two adjacent base
loci, with $1\le f=\dim V_t<d$. At a general incidence point the
entire higher-kernel base ideal satisfies
\[
 \mathfrak b_t\subset I_{V_t}^h,\qquad
 h=\lceil k\delta/2\rceil,\qquad k/h\le2/\delta.
\]
Consequently, on a resolution $V_t^*$,
\begin{equation}\label{eq:full-kernel-bounds}
 \begin{split}
 0<P_W^fV_t&\le(2/\delta)^{d-f}P_W^d,\\
 K_{V_t^*}P_W^{f-1}&\le(K_W+cP_W)|_{V_t}P_W^{f-1},
 \qquad 0<c\le2(d-f)/\delta.
 \end{split}
\end{equation}
\end{lemma}
\begin{proof}
On the constant-rank open, the full kernels are vector bundles
with nested proper base loci. Track a component through the mark
across the $d+1$ levels. Its dimensions lie in
$\{1,\ldots,d-1\}$, so two adjacent tracked components coincide.
After a finite extension of the parameter field, spread them out
to an integral family whose marked incidence dominates $W$.
The adjacent orders satisfy
\[
 \lceil k\tau_{i+1}\rceil-\lceil k\tau_i\rceil
 \ge k\delta-1\ge h-1.
\]
Lemma~\ref{lem:jet-normal-containment} gives the ordinary-power
containment for the entire higher base ideal. Apply
Lemmas~\ref{lem:hd-base-degree} and \ref{lem:canonical-centre}
for the degree and canonical bounds. The construction follows
\cite[Lemmas~4.5--4.6]{NS}.
\end{proof}

\begin{proposition}\label{prop:hd-moving-centre}
Let $P$ be an ample rational Cartier divisor on $X$, and put
$\rho=(P^3)^{1/3}$. Suppose that, for some positive integer $k_*$
with $k_*P$ Cartier, a Zariski dense set of points $x$ admits
$E_x\in|k_*P|$ with $\mult_x(E_x)>4k_*\rho$.
Then there is a covering family of integral subvarieties $V$ of
one fixed dimension $d\in\{1,2\}$ such that
\[
 0<P^d\cdot V\le 3^{3-d}(P^3)^{d/3}.
\]
Moreover, on a resolution $Y\to V$, with $P_Y$ the pullback,
\begin{equation}\label{eq:moving-canonical}
 K_YP_Y^{d-1}\le\frac{3(3-d)}{\rho}P_Y^d.
\end{equation}
The marked incidence dominates $X$.
\end{proposition}
\begin{proof}
Lemma~\ref{lem:hd-dense-multiplicity} makes the multiplicity
condition valid at every point. Taking powers of the defining
sections allows $k$ to be arbitrarily large and divisible by $k_*$.
Put
\[
 \delta=\frac{2\rho}{3},\qquad
 \tau_i=\frac32\rho+i\delta\quad(0\le i\le3),\qquad
 q_i=\lceil k\tau_i\rceil.
\]
Choose $k$ so that $k\delta\ge4$ and $k\rho/2>1$.
All full kernel spaces
\[
 V_i(x)=H^0(X,\OO_X(kP)\otimes\mathfrak m_x^{q_i})
\]
are nonzero, because $\tau_3=7\rho/2<4\rho$.
Write $B_i(x)$ for their base loci. They are proper and nested.
The point $x$ cannot be isolated in $B_0(x)$: otherwise
Lemma~\ref{lem:hd-base-degree} would give
$q_0^3\le k^3P^3$, whereas $q_0\ge3k\rho/2$.

Apply Lemma~\ref{lem:full-kernels} with $W=X$ and these four
levels. It gives a common component $V_t$ of dimension
$d\in\{1,2\}$ and a covering incidence, together with
ordinary-power containment for the entire higher base ideal,
where $h=\lceil k\delta/2\rceil$.

The higher base ideal is primary there, since $V_t$ is its component.
Lemma~\ref{lem:hd-base-degree}, applied to $H=kP$, gives
\[
 h^{3-d}k^d P^d\cdot V_t\le k^3P^3,
 \qquad
 P^d\cdot V_t\le(k/h)^{3-d}P^3
 \le(3/\rho)^{3-d}P^3.
\]
Lemma~\ref{lem:canonical-centre} gives
$K_YP_Y^{d-1}\le(3-d)(k/h)P_Y^d$, because $K_X\sim0$.
Since $k/h\le3/\rho$, this is \eqref{eq:moving-canonical}.
Ampleness gives positivity, and the incidence is dominant.
\end{proof}

\subsection{The scalar bound and the movable class}\label{sec:scalar}

\begin{lemma}\label{lem:bounded-curves}
Let $Y$ be an auxiliary smooth projective variety of dimension
$f\in\{1,2,3\}$. Let $B_0$ be a reduced simple normal crossing
divisor or zero, $L$ a nef Cartier divisor of full nef dimension,
and $P_Y$ a nef rational Cartier divisor. If
\[
 0<P_Y^f\le C_0,\quad P_Y-rL\text{ is nef},\quad
 K_Y+B_0+L\text{ is big},
\]
and
\[
 (K_Y+B_0+L)P_Y^{f-1}\le C_1P_Y^f,
\]
then
\[
 r\le(m_fC_1)^{f-1}C_0,
\]
where $m_f$ depends only on $f$.
\end{lemma}
\begin{proof}
Birkar--Zhang \cite[Theorem~1.3]{BZ}, with boundary coefficients
in $\{0,1\}$ and Cartier index one for $L$, makes
$|m_f(K_Y+B_0+L)|$ birational. Resolve its base ideal and let
$H$ be the free moving part. Then $H$ is nef and big and
\[
 HP_Y^{f-1}\le m_fC_1P_Y^f.
\]
Khovanskii--Teissier gives
\[
 P_YH^{f-1}\le P_Y^f
 \left(\frac{HP_Y^{f-1}}{P_Y^f}\right)^{f-1}
 \le(m_fC_1)^{f-1}C_0.
\]
General hyperplanes through the image of a very general point
under the birational morphism defined by $H$ cut an effective
curve cycle containing a curve $C$ through that point.
Full nef dimension and integrality give $LC\ge1$.
Thus $r\le P_YC\le P_YH^{f-1}$. For $f=1$ take $C=Y$.
This is \cite[Lemma~4.1]{NS}; in particular its constants are
independent of the varying auxiliary variety.
\end{proof}

\begin{proposition}\label{prop:scalar-bound}
Assume $\nef(X,D)=3$ and $D^3=0$. With
$P_n=r_n(D+A/n)$ as in Lemma~\ref{lem:hd-normalization}, put
$\rho_n=(P_n^3)^{1/3}$. For all sufficiently large $n$, and at a very general point $x$, every
member $E\in|kP_n|$, for every positive integer $k$ with $kP_n$
Cartier, satisfies
\[
 \mult_x(E)\le4k\rho_n.
\]
\end{proposition}
\begin{proof}
Suppose violations persist as $n\to\infty$. For fixed $n,k$, the
high-multiplicity condition is closed in the incidence space, so
a violation at a sufficiently very general point yields a
dense set of violations. Use
Proposition~\ref{prop:hd-moving-centre}, together with its canonical
bound \eqref{eq:moving-canonical}, and write $Y\to V$ for the
resolution of a general member. For $f=\dim V\in\{1,2\}$ it gives
\[
 0<P_Y^f\le3^{3-f}\rho_n^f,\qquad
 K_YP_Y^{f-1}\le\frac{3(3-f)}{\rho_n}P_Y^f.
\]
The volume window $P_n^3\to\epsilon^3>0$ makes both bounds
uniform in $n$. Let $L$ be the pullback of $D$; then
$P_Y-r_nL$ is nef.

Choose the incidence point above the very general full-nef locus
and outside the images of all families of rational curves.
The latter exclusion is possible because $K_X\sim0$ makes $X$
non-uniruled. Since $Y\to V\subset X$ is generically finite,
$Y$ is non-uniruled and $L$ has full nef dimension; this is
the very-general-images argument of \cite[Lemma~4.2]{NS}.
By \cite{BDPP}, $K_Y$ is pseudo-effective, and
Lemma~\ref{lem:lower-dimensional} makes $K_Y+L$ big.
Also $LP_Y^{f-1}\le P_Y^f/r_n$.
Lemma~\ref{lem:bounded-curves} now bounds $r_n$ uniformly,
contrary to $r_n\to\infty$.
There are only countably many $n,k$ and jet orders, so one
very general locus handles all the stated section spaces.
\end{proof}

\begin{proposition}\label{prop:movable-from-scalar}
Assume the scalar bound at a very general point $x$, and fix
large $n$ with $\rho_n<2\epsilon$ and $r_n>Q$.
On the blowup $\pi:\widehat X\to X$ at $x$ there is an actual
strongly movable class $\gamma=f_*(A_1A_2)$ such that
\[
 J\gamma>0,\qquad
 \pi^*(P+\lambda D)\gamma<40\epsilon J\gamma
 \quad(0\le\lambda\le Q).
\]
Here $J$ is exceptional, $f$ is one smooth birational model,
and the $A_i$ are ample integral Cartier divisors.
\end{proposition}
\begin{proof}
The class $2\pi^*P-40\epsilon J$ is not pseudo-effective.
Otherwise its convex combination with the big class $2\pi^*P$
would make $2\pi^*P-cJ$ big for rational
$8\rho_n<c<40\epsilon$. An effective divisor in a divisible
multiple of this class would push down to a member of $|2kP|$
with multiplicity exceeding the bound in
Proposition~\ref{prop:scalar-bound}.
Movable-curve duality \cite{BDPP} gives a generating class
$\gamma=f_*(A_1A_2)$ with
\[
 2\pi^*P\gamma<40\epsilon J\gamma.
\]
Strict negativity permits rational approximation of the ample
classes and clearing denominators. Nefness gives $J\gamma>0$.
For $0\le\lambda\le Q$ the divisor
\[
 2P-(P+\lambda D)=(r_n-\lambda)D+r_nA/n
\]
is nef, which gives the required inequalities.
Fix this actual model and complete intersection before
reducing characteristic.
\end{proof}

\begin{proof}[Proof of Lemma~\ref{lem:key-scalar}]
Fix $Q$ and take $n$ sufficiently large that
Proposition~\ref{prop:scalar-bound} applies, $\rho_n<2\epsilon$,
and $r_n>Q$. Proposition~\ref{prop:movable-from-scalar}
gives the required class and inequalities.
\end{proof}

\section{Proof of Lemma~\ref{lem:key-jets}}\label{sec:large-jets}

\subsection{The projective bundle and adjoint bigness}

For $P=P_n$ and $Q>0$, write $R=R_n$. On a slice over one fixed
base coordinate, put
\[
 Z_{\rm sl}=\mathbb P(\OO_X\oplus\OO_X(D)),\qquad
 R_{\rm sl}=P+Q\xi.
\]
We use $\pi_*\OO_Z(a)=\operatorname{Sym}^a
(\OO_X(D)_1\oplus\OO_X(D)_2)$; thus $\dim Z=7$ and
$\dim Z_{\rm sl}=4$.

\begin{lemma}\label{lem:bundle-volume}
With $P_a=P+aD$, the projective-bundle formulas give
\[
 R^7=140\int_0^Q P_a^3P_{Q-a}^3\,da,\qquad
 R_{\rm sl}^4=4\int_0^Q P_a^3\,da.
\]
For fixed $Q$, these tend to $140Q\epsilon^6$ and
$4Q\epsilon^3$, respectively, as $n\to\infty$. Also
\[
 K_Z\sim D_1+D_2-2\xi,\qquad K_{Z_{\rm sl}}\sim D-2\xi.
\]
\end{lemma}
\begin{proof}
For a split rank-two bundle,
$\pi_*(\xi^j)=\sum_{a+b=j-1}D_1^aD_2^b$ for $j\ge1$.
Expand $R^7$ and integrate the binomial terms to obtain the first
identity; its coefficient is $7!/(3!3!)=140$.
The analogous one-factor expansion gives the second identity.
Since $D^3=0$, $P^3\to\epsilon^3$, and $r_n\to\infty$,
$(P+aD)^3/P^3\to1$ uniformly for $0\le a\le Q$.
The two limits follow. The canonical formulas are the rank-two
projective-bundle formula with $K_X\sim0$.
\end{proof}

\begin{lemma}\label{lem:semiample-trivial-rays}
Let $U,N$ be rational Cartier divisors on a normal projective
variety, with $U$ semiample and $dU$ Cartier and free for an
integer $d>0$. Let $0<s_0<s$ satisfy
\[
 (1-s_0/s)/d>12s_0/s.
\]
If $C$ generates a $(U+s_0N)$-negative extremal ray and
$(U+sN)C\ge-12$, then $UC=0$. A contraction or flip of such
a ray preserves the semiampleness of $U$.
\end{lemma}
\begin{proof}
If $UC>0$, integrality gives $UC\ge1/d$, and
\[
 (U+s_0N)C=(1-s_0/s)UC+(s_0/s)(U+sN)C
 \ge(1-s_0/s)/d-12s_0/s>0,
\]
a contradiction. Thus $UC=0$ on every contracted ray.
The free line bundle $dU$ descends through a contraction and
pulls back through a flip, preserving semiampleness.
\end{proof}

\begin{lemma}\label{lem:positive-boundary-model}
Let $G\ge0$ be rational with $(X,G)$ klt and $\kappa(G)>0$.
There exist $s_0>0$ rational, a normal projective model $Y$, and
a common resolution $p:W\to X$, $q:W\to Y$, with
$U=K_Y+G_Y$ semiample, $Q_0=U+s_0N$ nef, and
\begin{equation}\label{eq:positive-boundary-comparison}
 p^*G=q^*U+E_0,\qquad
 p^*(G+s_0D)=q^*Q_0+E_1,
 \qquad E_0,E_1\ge0\text{ $q$-exceptional}.
\end{equation}
Moreover, $\kappa(U)=\kappa(G)$, $\nef(Y,Q_0)=3$, and $N$
is the trace of the nef b-divisor determined by $D$.
\end{lemma}
\begin{proof}
We spell out \cite[Proposition~3.3]{NS} in the only threefold
instance required below. Run the ordinary $G$-MMP to a good
model; ordinary threefold abundance is \cite{KollarEtAl}.
Its adjoint $U=K_{Y_0}+G_{Y_0}$ is semiample. Let $N_0$ be the
trace of the fixed nef b-divisor determined by $D$ on $X$.
For small $s>0$, this finite first program is also negative for
$G+sD$, and $(Y_0,G_{Y_0}+s\mathbf D)$ is generalized klt.


Choose $dU$ free and $s_0$ as in
Lemma~\ref{lem:semiample-trivial-rays}. Run a generalized
$(U+s_0N_0)$-program. The nef data are rational and NQC;
ordinary klt approximation gives its steps, and
\cite[Theorem~1.2]{ChenTsakanikas} gives termination.
The generalized length bound is $12$ in this threefold setting
\cite[Lemma~3.1]{NS}. Lemma~\ref{lem:semiample-trivial-rays}
makes every step $U$-trivial.
On the final model $Y$, $U$ remains
semiample, $Q_0=U+s_0N$ is nef, and a common resolution gives the two displayed comparisons.

Moreover $Q_0$ has full nef dimension. Otherwise a very general
compact nef-reduction fibre has a curve through its very general
point avoiding $q(\operatorname{Exc}q)$, which has codimension
at least two. Nefness of $U$, pseudo-effectivity of $N$, and
generic nonnegativity of pseudo-effective degrees give $UC=NC=0$.
Lifting this curve in both comparisons gives $D$-degree zero
through a very general point of $X$, a contradiction.
\end{proof}

\begin{lemma}\label{lem:positive-boundary}
If $G\ge0$ is rational, $(X,G)$ is klt, and $\kappa(G)>0$,
then $G+cD$ is big for every rational $c>0$.
\end{lemma}
\begin{proof}
Use the model in Lemma~\ref{lem:positive-boundary-model}.
The contraction $h:Y\to S$ defined by $U$ has
$\dim S=\kappa(G)>0$ and $U\sim_\Q h^*A_S$ with $A_S$ ample.
On a very general positive-dimensional fibre $F$,
\[
 K_F+G_Y|_F\sim_\Q0,
 \qquad Q_0|_F=s_0N|_F\text{ is nef of full nef dimension}.
\]
Its dimension is at most two, so Lemma~\ref{lem:lower-dimensional}
makes $Q_0|_F$ big. For zero-dimensional fibres relative bigness
is automatic. Relative bigness follows on the generic fibre: for a fixed ample
$H$, one of the countably many closed section loci for
$(\ell Q_0-H)|_{F_s}$ contains a very general point, hence is the
whole shrunken flat base. Adding a sufficiently large multiple
$bU$ now makes $Q_0+bU$ big globally. The comparisons imply
that $(1+b)G+s_0D$ is big. Finally, for
$0<\lambda<\min\{1/(1+b),c/s_0\}$,
\[
 G+cD=\lambda((1+b)G+s_0D)
 +(1-\lambda(1+b))G+(c-\lambda s_0)D
\]
is big.
\end{proof}

\begin{lemma}\label{lem:multisection-axes}
Assume $\nef(X,D)=3$. Let $F\subset
\mathbb P(\OO_X\oplus\OO_X(D))$ be an integral multisection
of degree $e>0$ meeting the torus open. Its two axis intersections
push down to effective integral divisors $G_1,G_2$ on $X$ satisfying
\begin{equation}\label{eq:axis-divisors}
 G_1-G_2\sim\pm eD,\qquad
 QG_iP^2\le R_{\rm sl}^3F.
\end{equation}
For $G=G_1+G_2$ and every sufficiently small rational $b>0$,
$(X,bG)$ is klt, $\kappa(bG)>0$, and $bG+D$ is big.
\end{lemma}
\begin{proof}
The norm of the ratio of the two fibre coordinates gives the
linear equivalence. If $S$ is either axis, $R_{\rm sl}-QS$
and $R_{\rm sl}-P$ are nef. Intersecting these classes against
the effective cycle $S\cdot F$ gives the degree inequalities.

If $\kappa(bG)=0$, ordinary threefold abundance gives a good
$(K_X+bG)$-model $Y$ with $K_Y+bG_Y\sim_\Q0$.
The transform of $K_X$ is linearly trivial, so $G_Y=0$:
every component of $G$ has been contracted. On a common
resolution $a:U\to X$, $c:U\to Y$, the signed relation in
\eqref{eq:axis-divisors} expresses $a^*D\sim_\Q E$, with
$E$ $c$-exceptional. Divisors exceptional over $X$ are also
exceptional over $Y$, since this MMP extracts none.
Since $a^*D$ is nef, negativity gives $E\le0$.
Its class is pseudo-effective, so intersection with an ample
square forces $E=0$. This contradicts full nef dimension.
Thus $\kappa(bG)>0$, and Lemma~\ref{lem:positive-boundary}
gives bigness. The axis construction and negativity argument
are dimension-independent; only the ordinary good-model input
is being specialized to threefolds.
\end{proof}

\begin{lemma}\label{lem:fixed-graphs}
Let $T_1,T_2$ be fixed non-uniruled integral projective varieties
over $\C$ of the same dimension. If an integral algebraic family
of graphs of birational maps $T_1\dashrightarrow T_2$ dominates
$T_1\times T_2$, both varieties are birational to an abelian variety.
\end{lemma}
\begin{proof}
Identify the two varieties by one map in the family.
Hanamura's theorem, in the form
\cite[Theorem~3.8]{Blanc}, supplies a smooth projective birational
model $Y$ with $\operatorname{Bir}(Y)_{\rm red}$ a group scheme
locally of finite type. Its identity component is
$\operatorname{Aut}^0(Y)$, an abelian variety acting regularly.

Conjugate the family to $Y$, take the closure of its relative
graphs, and shrink the integral parameter space so that these
graphs are flat and geometrically integral, with both projections
birational. This shrinking is necessary: graph flatness is not
preserved automatically by birational conjugation
\cite[Remark~3.6]{Blanc}. The flat-graph representability theorem
\cite[Proposition~3.7]{Blanc} then gives a morphism from the
reduced parameter variety to $\operatorname{Bir}(Y)_{\rm red}$.
Its connected image lies in one translate of the identity component.
After translating by a fixed member, product dominance implies
that a general point has a dense orbit under that abelian group.
Such an orbit is the quotient by its closed stabilizer, hence an
abelian variety. It is therefore a birational model of $Y$.
No canonical nonvanishing theorem on $T_i$ is used.
\end{proof}

\begin{lemma}\label{lem:abelian-adjoint}
Let $p:Y\to A_0$ be a birational morphism from a smooth projective
variety to an abelian variety. If $L$ is nef Cartier of full nef
dimension on $Y$, then $K_Y+L$ is big.
\end{lemma}
\begin{proof}
The numerical pushforward $N=p_*L$ is pseudo-effective on $A_0$,
thus nef. If it were not big, the kernel of the semipositive
Hermitian form of its integral class would define a
positive-dimensional abelian subvariety on whose translates $N$
is numerically trivial. The image of the exceptional locus of $p$
has codimension at least two. A general translate meets that image
in codimension at least two; complete-intersection curves in the
translate can be chosen through a very general point avoiding it.
Their strict lifts have $L$-degree zero, a contradiction. Thus $N$
is big.

Negativity gives $p^*N-L=E_0\ge0$, with $E_0$ exceptional.
Since $A_0$ is smooth and $K_{A_0}\sim0$, $K_Y\sim E_K\ge0$,
with positive coefficient along every $p$-exceptional divisor.
Choose $0<u<1$ rational with $E_K-uE_0\ge0$. Then
\[
 K_Y+uL\sim_\Q up^*N+(E_K-uE_0)
\]
is big. Adding $(1-u)L$ proves the assertion. This lemma has no
dimension restriction.
\end{proof}

\subsection{Base components and slice exclusions}

Fix $\delta=8\epsilon$, $\tau_i=9+i\delta$ for $0\le i\le7$,
and choose $Q$ so that $\tau_7^7<70Q\epsilon^6$.
For all sufficiently large $n$, Lemma~\ref{lem:bundle-volume} gives
\[
 \epsilon^3/2<P^3<2\epsilon^3,\qquad
 \tau_7^7<R^7<280Q\epsilon^6,\qquad R_{\rm sl}^4<8Q\epsilon^3.
\]
In particular $(2/\delta)^3\,8\epsilon^3=1/8<1$.
For the following lemmas, these choices remain fixed, and
$C_0=(2/\delta)^4\,280Q\epsilon^6$.

\begin{lemma}\label{lem:sevenfold-centres}
For all sufficiently large $n$, if $\varepsilon(R;z)\le8$ at every
point of the torus open, there is a covering family of proper
integral $V\subset Z$, with $f=\dim V$, such that
\begin{equation}\label{eq:sevenfold-centres}
 0<R^fV\le(2/\delta)^{7-f}R^7,\qquad
 K_{V^*}R^{f-1}\le\frac{2(7-f)}{\delta}R^fV.
\end{equation}
For a general fibre component $F$ of either projection, if
$1\le s=\dim F\le3$, the actual restricted system gives
\begin{equation}\label{eq:slice-centres}
 \begin{split}
 0<R_{\rm sl}^sF&\le(2/\delta)^{4-s}R_{\rm sl}^4,\\
 K_{F^*}R_{\rm sl}^{s-1}
 &\le(K_{Z_{\rm sl}}+c_sR_{\rm sl})|_F R_{\rm sl}^{s-1},
 \quad c_s\le2(4-s)/\delta.
 \end{split}
\end{equation}
Both images of a general marked incidence point are very general on $X$.
\end{lemma}
\begin{proof}
Asymptotic Riemann--Roch and $\tau_7^7<R^7$ make the eight full
kernels nonzero for large divisible $k$. A curve through a very
general $z$ with $R$-degree divided by multiplicity less than
$\tau_0$ is contained in the lowest base locus, by the intersection
multiplicity inequality. Hence that base locus is positive-dimensional.
Lemma~\ref{lem:full-kernels} gives the covering family and the first
bounds. Since $K_Z$ is minus the axes and $V$ meets their complement,
the ambient canonical contribution is nonpositive.
Apply Lemma~\ref{lem:restricted-components} to either projection
over its actual image, and then apply
Lemmas~\ref{lem:hd-base-degree} and \ref{lem:canonical-centre}
in the smooth fourfold slice. This gives the fibre bounds.
\end{proof}

\begin{lemma}\label{lem:slice-small-image}
For all sufficiently large $n$, a fibre $F$ as in
\eqref{eq:slice-centres} cannot be generically finite onto an
image $W_0\subset X$ of dimension less than three.
\end{lemma}
\begin{proof}
If $F\to W_0$ is generically finite and $\dim W_0<3$, the
pullback $L$ of $D$ is nef Cartier of full nef dimension on $F^*$,
and $K_{F^*}$ is pseudo-effective by the very-general-images
argument. Lemma~\ref{lem:lower-dimensional} gives $K_{F^*}+L$
big. The testing class $P_Y=R_{\rm sl}|_{F^*}$ satisfies
$P_Y-r_nL$ nef, and \eqref{eq:slice-centres} gives uniform volume
and canonical ratios. Lemma~\ref{lem:bounded-curves} bounds $r_n$,
a contradiction.
\end{proof}

\begin{lemma}\label{lem:saturated-slices}
For all sufficiently large $n$, a fibre $F$ as in
\eqref{eq:slice-centres} cannot be the full projective-line bundle
over its image $W_0\subset X$, with $\dim W_0\le2$.
\end{lemma}
\begin{proof}
If $F\to W_0$ has one-dimensional fibres, $F$ is the entire
projective-line bundle over $W_0$. Put $w=\dim W_0=s-1$.
For $w=0$, \eqref{eq:slice-centres} gives
\[
 Q=R_{\rm sl}F\le(2/\delta)^3R_{\rm sl}^4<Q.
\]
For $w=1,2$, resolve $W_0$ and take its induced bundle
$p_0:F^*\to W_0^*$. Cancelling the relative canonical term
$-2\xi+p_0^*D$, which equals the ambient canonical restriction,
gives
\[
 p_0^*K_{W_0^*}R_{\rm sl}^w\le c_sR_{\rm sl}^{w+1}F.
\]
Its term linear in $\xi$ is $wQK_{W_0^*}P^{w-1}$; all other
terms are nonnegative because $K_{W_0^*}$ is pseudo-effective
and $P,D$ are nef. For $r_n\ge Q$, expansion gives
\begin{equation}\label{eq:saturated-estimate}
 \begin{split}
 (w+1)QP^wW_0&\le R_{\rm sl}^{w+1}F\le7QP^wW_0,\\
 K_{W_0^*}P^{w-1}&\le(7c_s/w)P^wW_0.
 \end{split}
\end{equation}
The first lower bound and \eqref{eq:slice-centres} bound $P^wW_0$
uniformly. Applying Lemmas~\ref{lem:lower-dimensional} and
\ref{lem:bounded-curves} on $W_0^*$ is again contradictory.
\end{proof}

\begin{lemma}\label{lem:multisection-slices}
For all sufficiently large $n$, a fibre $F$ as in
\eqref{eq:slice-centres} cannot be a multisection over $X$.
\end{lemma}
\begin{proof}
Lemma~\ref{lem:multisection-axes} gives $G=G_1+G_2$, $bG+D$ big,
and $GP^2\le2R_{\rm sl}^3F/Q$. On a log resolution $a:Y\to X$,
let $B_0$ be the reduced strict support of $G$ together with every
exceptional divisor. Decrease $b$ so that
\[
 K_Y+B_0+a^*D\sim_\Q a^*(bG+D)+E_b,\qquad E_b\ge0.
\]
This adjoint is big. Testing with $P_Y=a^*P$, all exceptional
divisors contribute zero, and
\[
 (K_Y+B_0+a^*D)P_Y^2\le GP^2+DP^2\le C,
 \qquad DP^2\le P^3/r_n.
\]
The positive lower bound on $P^3$ gives a uniform adjoint ratio.
Lemma~\ref{lem:bounded-curves}, with this reduced boundary,
excludes the case.
\end{proof}

\begin{lemma}\label{lem:finite-correspondences}
For all sufficiently large $n$, both projections of a general
$V$ in Lemma~\ref{lem:sevenfold-centres} are generically finite
and dominant. On a simultaneous resolution $Y$, with projections
$g_i:Y\to X$ and $P_Y=R|_Y$,
\begin{equation}\label{eq:correspondence-estimates}
 0<P_Y^3\le C_0,\qquad K_YP_Y^2\le(8/\delta)P_Y^3,
 \qquad P_Y-r_ng_i^*D\text{ is nef}.
\end{equation}
Moreover, $\deg g_i\le2C_0/\epsilon^3$.
\end{lemma}
\begin{proof}
Lemmas~\ref{lem:slice-small-image}, \ref{lem:saturated-slices},
and \ref{lem:multisection-slices} exclude fibre dimensions one,
two, and three. A four-dimensional fibre makes $V$ the full
inverse image of a proper $W_0\subset X$, with $\dim W_0\le2$;
the other projection then has fibre dimension $\dim W_0+1$, already
excluded. Thus both projections are generically finite and $f\le3$.
For $f<3$, \eqref{eq:sevenfold-centres} and
Lemmas~\ref{lem:lower-dimensional} and \ref{lem:bounded-curves}
give a uniform bound on $r_n$, a contradiction. Hence $f=3$.
The displayed estimates follow from \eqref{eq:sevenfold-centres}
and nefness of $R-r_nD_i$.
Since $R-P_i$ is nef, $\deg g_i\,P^3\le P_Y^3\le C_0$.
\end{proof}


\subsection{Finite covers and the final exclusion}

\begin{lemma}\label{lem:branch-surfaces}
For all sufficiently large $n$, no relative Jacobian component
in the family of Lemma~\ref{lem:finite-correspondences} whose
general image is a divisor has images sweeping $X$.
\end{lemma}
\begin{proof}
Suppose that the images of one branch-divisor component sweep
$X$, and let $J\subset X$ be a general such image. Pulling back
a nowhere-vanishing three-form from $X$ gives an effective
canonical divisor on $Y$, containing a ramification divisor
over $J$ with positive integral coefficient. Thus
\[
 0<d_J:=P^2J\le K_YP_Y^2\le(8/\delta)C_0.
\]
Generic numerical subadjunction for $(X,J)$ and Hodge index give
on a resolution $J^*$
\begin{equation}\label{eq:branch-surface}
 K_{J^*}P\le J^2P
 \le\frac{(JP^2)^2}{P^3}\le C_2d_J,
 \qquad C_2=16C_0/(\delta\epsilon^3).
\end{equation}
Here $(X,J)$ is lc at $\eta_J$; global lc is unnecessary.
Hodge index applies to $J$ without assuming $J$ nef.
Since the images sweep, $J^*$ is non-uniruled and the pulled-back
$D$ has full nef dimension. The testing volume is $d_J$, bounded
above, its canonical ratio is bounded by $C_2$, and $P-r_nD$ is
nef. Lemmas~\ref{lem:lower-dimensional} and
\ref{lem:bounded-curves} give a uniform bound for $r_n$.
Hence sweeping branch divisors are impossible for large $n$.
\end{proof}

\begin{lemma}\label{lem:finite-stein-covers}
Fix sufficiently large $n$ and a smooth integral parameter family
as in Lemma~\ref{lem:finite-correspondences}. The normal Stein
covers of either projection belong to a finite list.
\end{lemma}
\begin{proof}
For this fixed $n$, retain the relative Jacobian components whose
general fibre images are divisors on $X$. Remove their finitely
many proper image closures, obtaining a fixed smooth open $U$. Every finite normal Stein cover of a
general $g_i$ is unramified at the codimension-one points of $U$;
purity makes it finite \'etale over $U$. The fundamental group of
$U$ is finitely generated, so it has finitely many subgroups of
the bounded indices in question. Normal covers of $X$ are
determined by their function fields. There are consequently only
finitely many possible Stein covers in each slot, as in
\cite[Lemma~7.5]{LA}. Degree bounds alone would not imply this.
\end{proof}

\begin{lemma}\label{lem:bounded-correspondence-exclusion}
A family as in Lemma~\ref{lem:finite-correspondences}, with
finitely many Stein covers in each slot, bounds $r_n$ by a constant
depending only on $\epsilon,Q,\delta$. Consequently such families
cannot occur for arbitrarily large $n$.
\end{lemma}
\begin{proof}
Each $Y$ is birational to both of its Stein covers. Fix a pair
$T_1,T_2$ occurring among the original members satisfying
\eqref{eq:correspondence-estimates}. The projected incidence
dominates $X\times X$. The possible graph families for the
finitely many pairs belong to countably many Hilbert schemes;
stratify for flat integral fibres and birational projections.
Uncountability of $\C$ implies that one such integral graph
family has evaluation dominant over $X\times X$. Its evaluation
is also dominant on $T_1\times T_2$: this product is integral,
finite over $X\times X$, and of the same dimension.

The covers are non-uniruled, since they are finite over $X$.
Lemma~\ref{lem:fixed-graphs} makes them birational to an abelian
variety $A_0$. Return to an \emph{original} correspondence for
this pair and resolve further to get a birational morphism
$p:Y\to A_0$. This preserves the numerical bounds
\eqref{eq:correspondence-estimates}; additional exceptional
canonical divisors have zero intersection with the pulled-back
testing square. The divisor $L=g_1^*D$ remains nef Cartier of
full nef dimension. Lemma~\ref{lem:abelian-adjoint} gives
$K_Y+L$ big, while
\[
 (K_Y+L)P_Y^2\le(8/\delta+1/r_n)P_Y^3.
\]
Lemma~\ref{lem:bounded-curves} now excludes this final case.
The cover lists and graph families may depend on $n$; their
sizes enter none of the uniform numerical bounds. Failures
therefore cannot persist along $n\to\infty$.
\end{proof}


\begin{proposition}\label{prop:sevenfold-jets}
There is a fixed integer $Q>0$ such that, for all sufficiently
large $n$, some point $z$ in the torus open satisfies
$\varepsilon(R;z)>8$.
\end{proposition}
\begin{proof}
Use the fixed choices above. If the assertion failed for arbitrarily
large $n$, Lemmas~\ref{lem:sevenfold-centres} and
\ref{lem:finite-correspondences} would produce bounded threefold
correspondences. Resolve and shrink the family, and spread out
the finitely many generic Jacobian components after a finite
parameter extension. Lemmas~\ref{lem:branch-surfaces} and
\ref{lem:finite-stein-covers} give finite lists of Stein covers.
Lemma~\ref{lem:bounded-correspondence-exclusion} contradicts
$r_n\to\infty$.
\end{proof}

\begin{proof}[Proof of Lemma~\ref{lem:key-jets}]
Choose $z$ as in Proposition~\ref{prop:sevenfold-jets}, and rational
$8<a<\varepsilon(R;z)$. On the blowup $b:\widetilde Z\to Z$,
$b^*R-aJ_z$ is ample. For large divisible $k_0$, Serre vanishing
makes restriction surjective modulo $\mathfrak m_z^{k_0a}$,
with $k_0a\ge8k_0+1$. This is the required evaluation surjection.
\end{proof}

\section{Proof of Lemma~\ref{lem:key-contradiction}}\label{sec:frobenius}

All characteristic-zero data are fixed before reduction. The
following estimates specialize \cite[Theorem~6.1]{NS}; Frobenius
is used only on smooth positive-characteristic fibres.

\subsection{Fixed data and generic rank}

\begin{lemma}\label{lem:cotangent-flag}
For the fixed ample rational divisor $P$, there is a smooth
complete intersection flag with successive classes $l_1P,l_2P$
and final curve $C$ such that
$\mu_{\min}(\Omega_X^1|_C)\ge0$.
\end{lemma}
\begin{proof}
Since $K_X\sim0$, generic cotangent semipositivity
\cite[Theorem~2.1]{CP} makes $\Omega_X^1$ semistable of slope
zero for this polarization. Mehta--Ramanathan restriction
\cite[Theorem~6.1]{MR} gives the assertion for sufficiently
large divisible $l_1,l_2$. Bertini makes the flag smooth and
integral. These results apply in arbitrary dimension; here
only two restrictions are needed.
\end{proof}

\begin{lemma}\label{lem:jet-products-hd}
If sections of a line bundle $L$ on $Z$ generate jets through
order $j$ at $z$, then products of $N$ such sections generate jets
through order $Nj$ for $L^N$.
\end{lemma}
\begin{proof}
Every monomial of degree at most $Nj$ is a product of $N$ monomials
of degrees at most $j$. Lift the latter to sections with those leading
terms and multiply. The product has the desired leading term and
only higher-order errors. In a basis ordered by total degree this
gives a triangular matrix with nonzero diagonal on the finite jet
quotient, hence surjectivity.
\end{proof}

\begin{proposition}\label{prop:frob-rank-hd}
Let $Y$ be a smooth positive-characteristic fibre of a model of
$X$, and use the same notation for the specialized divisor data, over an algebraically closed field of characteristic $p$.
On $\mathbb P(\OO_Y(D)_1\oplus\OO_Y(D)_2)$, suppose a fixed
finite system of sections of $k_0(P_1+P_2+Q\xi)$ generates jets
through order $8k_0$ at a point $z$ off the axes, where
$k_0P$ is Cartier. Choose $T$ so that all $T+aD$ for
$0\le a\le(k_0-1)Q$ have sections nonzero at the two base points.
Put $N_p=\lfloor p/k_0\rfloor$, $B_p=k_0N_pP+T$.
For $p\ge8k_0$, the two-slot evaluation map has generic rank
$p^6$.
\end{proposition}
\begin{proof}
Lemma~\ref{lem:jet-products-hd} gives order
$8k_0N_p$ for $N_pk_0R$. Put $d_p=(p-k_0N_p)Q$.
A weight-zero section of $d_p\xi+T_1+T_2$ nonzero at $z$ pads this
system to $pQ\xi+B_{p,1}+B_{p,2}$. Its multiplication is invertible
on the local jet quotients. Moreover
\[
 8k_0N_p\ge8(p-k_0+1)\ge7(p-1).
\]
In seven regular parameters, a monomial surviving modulo
$(u_1^p,\ldots,u_7^p)$ has degree at most $7(p-1)$.
Thus the padded sections surject onto this Frobenius quotient.

Let $F:Y\to Y'$ be relative Frobenius, let $D'$ be the twist, and
put $\mathcal E=F_*\OO_Y(B_p)$. Smoothness makes $F$ finite flat
of degree $p^3$, so $\operatorname{rk}\mathcal E=p^3$.
The tensor product $A_*$ of the two local Frobenius-fibre algebras
has dimension $p^6$. If $t$ is the fibre coordinate and $t(z)=c\ne0$,
the full quotient is $A_*[t]/(t^p-c^p)$, free over $A_*$ with
basis $1,t,\ldots,t^{p-1}$. Projection to the coefficient of $1$
keeps exactly weights $j=pa$, with $0\le a\le Q$.
Consequently the natural map
\[
 \begin{split}
 \bigoplus_{a+b=Q}&H^0(Y,B_p+paD)\otimes H^0(Y,B_p+pbD)\\
 &\quad\otimes\OO_{Y'\times Y'}(-aD'_1-bD'_2)
 \longrightarrow\mathcal E\boxtimes\mathcal E
 \end{split}
\]
spans $A_*$ at the base pair. The projection formula identifies
these summands with the evaluation map; the surviving fibre scalars
$c^{pa}$ are nonzero. The map therefore has generic rank $p^6$.
No reducedness assumption on $A_*$ is used.
\end{proof}

\begin{lemma}\label{lem:fixed-data-reduction}
Fix the jet system, padding sections, canonical trivialization,
flag, blowup, and complete-intersection witness over $\C$,
with $\mu_{\min}(\Omega_X^1|_C)\ge0$. They extend to a model whose smooth closed fibres occur in
arbitrarily large characteristics and preserve the evaluation
surjection, ampleness, intersection numbers, and semistability
of $\Omega_Y^1|_C$ of slope zero.
\end{lemma}
\begin{proof}
Choose the padding divisor $T$ and the finitely many nonvanishing
padding sections as in Proposition~\ref{prop:frob-rank-hd}.
Spread out the flag, canonical trivialization, jet sections,
blowup, and the actual complete intersection witness for
$\gamma$. Shrinking the model preserves smoothness, ampleness,
the finite jet surjection, and all intersection numbers.
Since $\deg(\Omega_X^1|_C)=K_XC=0$, the least-slope condition
makes this curve bundle semistable of slope zero. Openness of
semistability preserves it after further shrinking.
There remain closed fibres in arbitrarily large characteristics.
\end{proof}


On each smooth fibre let $F:Y\to Y'$ denote relative Frobenius.
For the following lemmas put
\[
 N_p=\lfloor p/k_0\rfloor,\quad B_p=k_0N_pP+T,\quad
 \mathcal E_p=F_*\OO_Y(B_p),\quad s_p=p^3,\quad R_p=p^6.
\]
All constants in $O(p^{-1})$ depend only on the fixed data.

\subsection{The diagonal rank}

\begin{lemma}\label{lem:truncated-slope-bound}
Let $V$ be a semistable vector bundle of rank three and degree
zero on a smooth projective curve of genus $g$ in characteristic
$p$. For $0\le j\le3(p-1)$, its truncated symmetric powers satisfy
\[
 \mu_{\max}(\mathcal A_j(V))\le3c_g,
\]
where $c_g$ depends only on $g$ and the rank, independently of $j,p$.
\end{lemma}
\begin{proof}
Langer's bound
\cite[Corollary~2.5]{Langer} gives
$L_{\min}(V)\ge-c_g/(p-1)$ with $c_g$ depending only on the
genus and rank. Put $u=3(p-1)$. The perfect complementary
degree pairing gives
\[
 \mathcal A_{u-j}(V)
 \simeq\mathcal A_j(V)^\vee\otimes(\det V)^{p-1}.
\]
Truncated symmetric powers are quotients of tensor powers;
their asymptotic least slopes are at least the corresponding
multiple of $L_{\min}(V)$. Since $\deg\det V=0$, the pairing
implies $\mu_{\max}(\mathcal A_j(V))\le3c_g$, uniformly in $j,p$.
\end{proof}

\begin{lemma}\label{lem:flag-section-bound}
Let $Y\supset S\supset C$ be a smooth complete intersection flag
cut by successive ample Cartier divisors $h_1,h_2$.
If $\mathcal G$ is locally free of rank $e$ and
$\mu_{\max}(\mathcal G|_C)\le M$, with $M\ge0$, then
\[
 h^0(Y,\mathcal G)\le e(M+1)
 \prod_{i=1}^2\left(1+\frac{M}{h_iC}\right).
\]
\end{lemma}
\begin{proof}
Successive restriction along the flag bounds sections by the sum
of the spaces on $C$ twisted by $-b_1h_1-b_2h_2$, with $b_i\ge0$.
Only pairs satisfying $b_1h_1C+b_2h_2C\le M$ contribute, since
negative maximal slope forbids sections. For each pair,
evaluation at $\lfloor M\rfloor+1$ points is injective and bounds
the dimension by $e(M+1)$.
There are at most $\prod_i(1+M/(h_iC))$ pairs. The restriction
sequences terminate because sufficiently negative ample twists
have no sections.
\end{proof}

\begin{lemma}\label{lem:diagonal-rank-bound}
On a smooth fibre of Lemma~\ref{lem:fixed-data-reduction}, put
$B_p=k_0\lfloor p/k_0\rfloor P+T$ and $R_p=p^6$.
If $(2P+QD)P^2<4P^3$ and $343P^3<1/4$, the weighted evaluation
map has rank less than $R_p/4$ on the diagonal of $Y'\times Y'$.
\end{lemma}
\begin{proof}
Let $\Delta_F=Y\times_{Y'}Y$, and put
\[
 \mathcal L_p=
 \OO_{\Delta_F}(B_{p,1}+B_{p,2}+pQD_1).
\]
On $\Delta_F$, the weighted section products all take values
in this same line bundle, because $pD_1=pD_2$ there.
Their value rank at any point of the diagonal $Y'\subset
Y'\times Y'$ is at most $h^0(\Delta_F,\mathcal L_p)$.
The canonical Frobenius filtration
\cite[Corollary~3.5]{KS} has graded pieces
\[
 \mathcal G_j=
 \OO_Y(2B_p+pQD)\otimes\mathcal A_j(\Omega_Y^1),
 \quad0\le j\le3(p-1),\qquad
 e_j=\operatorname{rk}\mathcal A_j,\quad \sum_j e_j=p^3,
\]
Here $\mathcal A_j$ is the truncated symmetric power of degree $j$,
locally the degree-$j$ part of
\[
 \OO[u_1,u_2,u_3]/(u_1^p,u_2^p,u_3^p).
\]


By Lemma~\ref{lem:truncated-slope-bound}, the truncated powers
on $C$ have uniformly bounded maximal slopes.
Set $\Lambda=l_1l_2$, so $[C]=\Lambda P^2$ and
$h_iC=l_i\Lambda P^3$. The differences
\[
 B_p-pP=T-(p-k_0N_p)P
\]
belong to a fixed finite list. The displayed intersection
hypothesis therefore gives one constant $c_0$, independent of
$j,p$, for which
\[
 \mu_{\max}(\mathcal G_j|_C)
 \le M_p:=7p\Lambda P^3+c_0.
\]

Lemma~\ref{lem:flag-section-bound} therefore gives
\[
 h^0(Y,\mathcal G_j)
 \le e_jp^3\bigl(343P^3+O(p^{-1})\bigr).
\]
Summing the graded pieces gives
\begin{equation}\label{eq:diagonal-rank}
 h^0(\Delta_F,\mathcal L_p)
 \le R_p\bigl(343P^3+O(p^{-1})\bigr)<R_p/4.
\end{equation}
The last inequality follows from
$343P^3<(14\epsilon)^3<1/4$ for large $p$.
There is no extra factor for the number of graded pieces:
their total rank is exactly $p^3$.
\end{proof}


\subsection{Determinant multiplicities}

\begin{lemma}\label{lem:determinant-line-bundles}
For the evaluation map in Proposition~\ref{prop:frob-rank-hd},
put $s_p=p^3$, $R_p=p^6$, and $\mathcal E_p=F_*\OO_Y(B_p)$.
There exist line bundles
\[
 \mathcal Q_i=(\det\mathcal E_p)^{s_p}\otimes\OO_{Y'}(m_iD'),
 \qquad m_i\in\mathbb Z_{\ge0},\quad m_1+m_2=QR_p,
\]
and a nonzero determinant section of $\mathcal Q_1\boxtimes\mathcal Q_2$.
Under the numerical identification with $Y'$, for $\lambda_i=m_i/R_p$,
\[
 \frac{c_1(\mathcal Q_i)}{R_p}=P+\lambda_iD+O(p^{-1}),
 \qquad0\le\lambda_i\le Q.
\]
\end{lemma}
\begin{proof}
Proposition~\ref{prop:frob-rank-hd} gives generic evaluation
rank $R_p$. Choose $R_p$ generically independent evaluation
columns of weights $(a,b)$, $a+b=Q$. Their determinant is
a nonzero section
\[
 \sigma\in H^0(Y'\times Y',\mathcal Q_1\boxtimes\mathcal Q_2),
 \qquad
 \mathcal Q_i=(\det\mathcal E_p)^{s_p}\otimes\OO_{Y'}(m_iD'),
\]
where $m_i\ge0$ and $m_1+m_2=QR_p$. Let $E_\sigma$ be the
effective divisor of $\sigma$, so
\[
 E_\sigma\in|\mathcal Q_1\boxtimes\mathcal Q_2|.
\]


The Frobenius formula for the first Chern class
\cite[Lemma~4.2]{Sun} gives, under the natural numerical
identification with the twisted fibre,
\[
 \frac{c_1(\mathcal Q_i)}{R_p}
 =\frac{B_p}{p}+\frac{p-1}{2p}K_Y+\lambda_iD
 =P+\lambda_iD+O(p^{-1}),\qquad
 \lambda_i=m_i/R_p\in[0,Q].
\]
\end{proof}

\begin{lemma}\label{lem:product-multiplicity-bound}
Let $L_1,L_2$ be line bundles on smooth projective varieties
$Y_1,Y_2$, and let $x_i\in Y_i$. If every $E_i\in|L_i|$ satisfies
$\mult_{x_i}(E_i)\le a_i$, then every $E\in|L_1\boxtimes L_2|$
satisfies
\[
 \mult_{(x_1,x_2)}(E)\le a_1+a_2.
\]
\end{lemma}
\begin{proof}
By K\"unneth the section space is the tensor product of the two
section spaces. Choose bases adapted to their multiplicity
filtrations. Their leading forms are independent in each degree,
and tensor products remain independent in each bidegree.
The largest possible nonzero initial degree is at most $a_1+a_2$.
\end{proof}

\begin{lemma}\label{lem:determinant-upper-bound}
Let $E_\sigma$ be the determinant divisor in
Lemma~\ref{lem:determinant-line-bundles}. The specialized movable
witness in Lemma~\ref{lem:fixed-data-reduction} gives
\begin{equation}\label{eq:det-upper}
 \mult_{(x',x')}(E_\sigma)
 \le R_p\bigl(80\epsilon+O(p^{-1})\bigr).
\end{equation}
\end{lemma}
\begin{proof}
Twist the specialized blowup and its fixed movable witness as well.
For every member $E_i\in|\mathcal Q_i|$, its proper transform
is effective. Intersecting it with the specialized class $\gamma$
gives
\[
 \frac{\mult_{x'}(E_i)}{R_p}\le40\epsilon+O(p^{-1}).
\]
The fixed complete intersection retains positivity on the
specialized model. The intersection inequality for effective divisors
therefore applies to every member of these linear systems on the fibre.

Apply Lemma~\ref{lem:product-multiplicity-bound} to the two
line bundles $\mathcal Q_i$.
\end{proof}

\begin{lemma}\label{lem:matrix-multiplicity-bound}
Let $M$ be an $R\times R$ matrix of regular functions near a
smooth point $u$, with $\det M\ne0$ and $\operatorname{rank}M(u)=r$.
Then
\[
 \mult_u(\operatorname{div}(\det M))\ge R-r.
\]
\end{lemma}
\begin{proof}
Constant invertible row operations make $R-r$ rows vanish at $u$.
Each determinant monomial contains one entry from each such row,
so belongs to $\mathfrak m_u^{R-r}$.
\end{proof}


\begin{proof}[Proof of Lemma~\ref{lem:key-contradiction}]
Since $P-r_nD$ is nef and $r_n>Q+1$,
\[
 (2P+QD)P^2\le(2+Q/r_n)P^3<4P^3,
 \qquad343P^3<(14\epsilon)^3<1/4.
\]
Apply Lemma~\ref{lem:fixed-data-reduction} and
Proposition~\ref{prop:frob-rank-hd} to the fixed data.
Lemma~\ref{lem:determinant-line-bundles} gives a nonzero
determinant divisor $E_\sigma$.
Its evaluation matrix has rank less than $R_p/4$ at $(x',x')$
by Lemma~\ref{lem:diagonal-rank-bound}. Thus
Lemma~\ref{lem:matrix-multiplicity-bound} gives
\begin{equation}\label{eq:det-lower}
 \mult_{(x',x')}(E_\sigma)>3R_p/4.
\end{equation}
Lemma~\ref{lem:determinant-upper-bound} gives the upper bound.
Taking $E_p=E_\sigma$ proves the assertion.
\end{proof}

\noindent\textit{Disclosure.} This AI-assisted report is intended for
joint learning, discussion, and checking. Its correctness has not been
independently verified. No claim of novelty, priority, or authorship
credit is made. Published and preprint inputs are attributed explicitly;
the scope of the source-dependent proof draft is stated above.

\begin{thebibliography}{99}
\bibitem{LA}
\emph{Log abundance in characteristic zero},
OpenAI mathematics collection, September 24, 2026,
Lemmas~7.1--7.5 (Tools for moving base components).
\href{https://github.com/openai/math/blob/main/preprints/Log-abundance-in-characteristic-zero-September-24-2026/paper.pdf}{Source manuscript}.
\bibitem{FH}
O.~Fujino and K.~Hashizume,
\emph{Existence of log canonical modifications and its applications},
Eur. J. Math. \textbf{9} (2023), Paper No.~13, Theorem~2.10.
\href{https://doi.org/10.1007/s40879-023-00598-0}{DOI}.
\bibitem{NS}
\emph{Numerical Semiampleness of Nef Adjoint Divisors},
OpenAI mathematics collection, October 3, 2026,
Proposition~3.3, Lemmas~4.1--4.7, Proposition~5.3,
Theorem~6.1, and Sections~8.4--8.5.
\href{https://github.com/openai/math/blob/main/preprints/Numerical-Semiampleness-of-Nef-Adjoint-Divisors-October-3-2026/paper.pdf}{Source manuscript}.
\bibitem{EKL}
L.~Ein, O.~K\"uchle, and R.~Lazarsfeld,
\emph{Local positivity of ample line bundles},
J. Differential Geom. \textbf{42} (1995), 193--219, Proposition~2.3.
\href{https://arxiv.org/pdf/alg-geom/9408003}{Preprint}.
\bibitem{LP}
V.~Lazi\'c and Th.~Peternell,
\emph{On Generalised Abundance, I},
Publ. Res. Inst. Math. Sci. \textbf{56} (2020), 353--389.
Numbering follows arXiv:1808.00438v2, especially Corollary~C
and Lemmas~2.14 and 3.1.
\href{https://arxiv.org/pdf/1808.00438}{Preprint}.
\bibitem{Ambro}
F.~Ambro,
\emph{The moduli $b$-divisor of an lc-trivial fibration},
Compos. Math. \textbf{141} (2005), 385--403, Theorem~0.2.
\href{https://doi.org/10.1112/S0010437X04001071}{DOI}.
\bibitem{FG}
O.~Fujino and Y.~Gongyo,
\emph{On the moduli $b$-divisors of lc-trivial fibrations},
Ann. Inst. Fourier \textbf{64} (2014), 1721--1735,
Definition~3.1 and Theorem~3.5.
\href{https://doi.org/10.5802/aif.2894}{DOI}.
\bibitem{BZ}
C.~Birkar and D.-Q.~Zhang,
\emph{Effectivity of Iitaka fibrations and pluricanonical
systems of polarized pairs},
Publ. Math. Inst. Hautes \'Etudes Sci. \textbf{123} (2016),
283--331, Theorem~1.3.
\href{https://doi.org/10.1007/s10240-016-0080-x}{DOI}.
\bibitem{KollarEtAl}
J.~Koll\'ar et al.,
\emph{Flips and Abundance for Algebraic Threefolds},
Ast\'erisque \textbf{211} (1992).
\href{https://www.numdam.org/item/AST_1992__211__1_0/}{Volume}.
\bibitem{ChenTsakanikas}
G.~Chen and N.~Tsakanikas,
\emph{On the termination of flips for log canonical generalized pairs},
Acta Math. Sin. (Engl. Ser.) \textbf{39} (2023), 967--994,
Theorem~1.2.
\href{https://arxiv.org/pdf/2011.02236}{Preprint}.
\bibitem{FujinoFG}
O.~Fujino,
\emph{Birkar--Cascini--Hacon--McKernan},
author's notes, version 1.02, November 28, 2009, Theorem~1.4.
\href{https://www.math.kyoto-u.ac.jp/~fujino/BCHM01.pdf}{Notes}.
\bibitem{HX}
C.~D.~Hacon and C.~Xu,
\emph{Existence of log canonical closures},
Invent. Math. \textbf{192} (2013), 161--195.
Numbering follows arXiv:1105.1169v3.
\href{https://doi.org/10.1007/s00222-012-0409-0}{DOI}.
\bibitem{Blanc}
J.~Blanc,
\emph{Algebraic structures of groups of birational transformations},
Proc. Sympos. Pure Math. \textbf{94} (2017), 17--30,
Theorems~3.8--3.9.
\href{https://algebra.dmi.unibas.ch/blanc/articles/structurebir.pdf}{Published offprint}.
\bibitem{BDPP}
S.~Boucksom, J.-P.~Demailly, M.~P\u{a}un, and Th.~Peternell,
\emph{The pseudo-effective cone of a compact K\"ahler manifold and
varieties of negative Kodaira dimension},
J. Algebraic Geom. \textbf{22} (2013), 201--248.
\href{https://doi.org/10.1090/S1056-3911-2012-00574-8}{DOI}.
\bibitem{CP}
F.~Campana and M.~P\u{a}un,
\emph{Orbifold generic semi-positivity: an application to families
of canonically polarized manifolds},
Ann. Inst. Fourier \textbf{65} (2015), 835--861, Theorem~2.1.
\href{https://doi.org/10.5802/aif.2945}{DOI}.
\bibitem{MR}
V.~B.~Mehta and A.~Ramanathan,
\emph{Semistable sheaves on projective varieties and their
restriction to curves},
Math. Ann. \textbf{258} (1982), 213--224, Theorem~6.1.
\href{https://doi.org/10.1007/BF01450677}{DOI}.
\bibitem{Langer}
A.~Langer,
\emph{Semistable sheaves in positive characteristic},
Ann. of Math. \textbf{159} (2004), 251--276, Corollary~2.5.
\href{https://doi.org/10.4007/annals.2004.159.251}{DOI}.
\bibitem{KS}
Y.~Kitadai and H.~Sumihiro,
\emph{Canonical filtrations and stability of direct images by
Frobenius morphisms},
Tohoku Math. J. \textbf{60} (2008), 287--301, Corollary~3.5.
\href{https://doi.org/10.2748/tmj/1215442876}{DOI}.
\bibitem{Sun}
X.~Sun,
\emph{Direct images of bundles under Frobenius morphism},
Invent. Math. \textbf{173} (2008), 427--447,
Theorem~3.7 and Lemma~4.2; numbering follows
arXiv:math/0611360v2.
\href{https://doi.org/10.1007/s00222-008-0125-y}{DOI}.
\end{thebibliography}
\end{document}
